
\subsection{Real-World Driving Benchmark}
\label{sec:method_driving}

To measure gaze accuracy during real driving, we associate
wearable gaze estimates with independently localised physical
targets and synchronised vehicle measurements
(Fig.~\ref{fig:framework}(a)).

\textbf{Vehicle and synchronisation.}
Our research vehicle is equipped with cameras, roof-mounted
LiDAR sensors and a GNSS/INS system. We use the front-left
camera, two 128-beam LiDARs and GNSS/INS measurements
for scene geometry and vehicle motion. Vehicle sensors share
a GNSS-disciplined clock via the Precision Time Protocol.
The Neon companion clock is aligned using the round-trip-based
Pupil Labs Time Echo protocol~\cite{pupil_time_echo},
enabling temporal association of gaze, images, LiDAR
and vehicle measurements.


\textbf{Reference generation.}
During driving, participants verbally announce identifiable
fixation targets, e.g., ``Fixating the truck on the right.''
(Fig.~\ref{fig:framework}c). Whisper transcribes the announcements,
and Hassil identifies the fixation event and object category.
An Ultralytics YOLO26 detector, trained on 3,827 annotated
images from the eye-tracker's scene-camera perspective,
proposes candidate targets within a generous region around
the estimated gaze position.
CoTracker3 point tracking~\cite{karaev2025cotracker3}
and RAFT optical flow~\cite{teed2020raft}
propagate target positions between detections and support
temporal consistency checks. Target proposals and tracking
are automated, with detection or tracking failures labelled
manually. The verified object centre defines the reference
independently of the measured gaze coordinates.

Knowing the announced target provides a reference for
identifying intervals with a stable gaze-to-target offset,
including smooth pursuit. Changes in this offset and eye
movement speed help distinguish target fixation from looking
away during driving. All intervals and target assignments
are manually verified in Pupil Player using gaze overlays
and infrared eye videos; 200 target assignments received a second
review. Inter-annotator agreement was not quantified. All finite,
metric-valid gaze and target samples within the verified intervals
are evaluated.



Participants were instructed to fixate the centres of
the announced targets. We select compact, clearly
identifiable objects such as traffic lights, licence
plates and road signs to minimise reference uncertainty.
Human gaze preferentially targets object
centres~\cite{kowler1995accuracy,nuthmann2010object}.
The median centre-to-corner extent is 0.64\dg{}
(P90: 1.30\dg{}), a geometric limit reached only
by fixating an extreme corner rather than the
instructed centre.




\textbf{Reference geometry and context.}
Verified target positions are converted into angular directions
using the calibrated scene-camera intrinsics and compared with
temporally aligned gaze estimates. Target distances are obtained from LiDAR measurements.
The vehicle-to-wearable association uses the sensor transformation
tree with a nominal head position; ordinary head movements
introduce uncertainty in the associated target distance.
For dense metric depth, we align Depth Anything predictions
with projected LiDAR measurements using robust per-frame
scale-and-shift fitting in inverse-depth space~\cite{marsal2025depth}. Each reference is associated with
vehicle motion and scene metadata. The resulting benchmark and
its driving conditions are characterised in
Section~\ref{sec:realworld}.


\textbf{Gaze suitability checks.}
To complement spatial accuracy, we assess whether gaze
measurements are suitable for downstream supervision.
A YOLO26 model trained on 548 annotated images from
three vehicles segments windshields, side windows and
mirrors to identify relevant viewing regions.
Image sharpness (Laplacian variance), optical flow and
IMU motion help detect motion-induced image degradation.
Pupil Labs blink events, eye openness and gaze dynamics
identify periods of potentially unreliable gaze estimation,
particularly during eyelid occlusion.
These manually validated indicators provide sample-level
quality flags for selecting usable gaze sequences.

